% !TEX root=jamc80.tex
\documentclass[document.tex]{subfiles}

\begin{document}
\section{\cite{jamc80} General equations}
Conservation of mass, momentum and energy: 

\textcolor{green}{ADD ALGEBRA COMPARISON SCHUBERT in final doc}
\begin{align}
    \pdv{\rho}{t} + \nabla \cdot (\rho\bm{u}) &= 0  \label{eq:3}\\
    \rho\left(\pdv{\bm{u}}{t} + \bm{u} \cdot \nabla \bm{u}\right) &= -\nabla P +\rho g + \pdv{x_j} \tau_{ij} \label{eq:4}\\
    \rho C_p \left\{\pdv{T}{t} - \frac{\alpha T}{\rho C_p}\pdv{P}{t} + \bm{u} \cdot [\nabla T - (\nabla T)_s]\right\} &= 
    \nabla \cdot (K\nabla T) + H + \tau_{ij}\pdv{u_i}{u_{ij}} \label{eq:5}
\end{align}

All terms and symbols used in these equations are shown in table \ref{table:symbols}. 
The (Newtonian) relation between the deviatoric stress tensor $\tau_{ij}$ and the velocity field is

\begin{equation}
\tau_{ij} = \eta \left( \frac{\partial u_i}{\partial x_j} + \frac{\partial u_j}{\partial x_i}
 - \frac{2}{3}\,\delta_{ij}\,\frac{\partial u_k}{\partial x_k} \right) 
 + \eta_2\,\delta_{ij}\,\frac{\partial u_k}{\partial x_k}.
\label{eq:8}
\end{equation}

Combining Eq.~\eqref{eq:8} with the momentum equation \eqref{eq:4} gives the general Navier-Stokes equations, 
but \textcite{jamc80} does not show these. Instead, they assume that $\frac{1}{Pr} \to 0$, which
they assume (\textcolor{green}{do they?}) is the case for the Earth's mantle, the inertial forces can be ignored. The resulting equations 
will be discussed in the next section, which covers the approximations made by \textcite{jamc80}. 

\section{\cite{jamc80} Approximations}
Problem statement: for the Boussinesq approximation to be valid, the depth range of convection $d$ must be 
small compared to the temperature scale height $H_T$ ($d/H_T \ll 1$). This is not the case for many systems of 
deep convection, amongst which the convection of planetary mantles. Physically, $d/H_T \approx O(1)$ means that 
the influence of viscous dissipation, variable adiabatic temperature- and density gradients must be maintained in the 
set of governing equations. 
\newline
The assumptions made for the \textbf{anelastic liquid approximation} are that $d=O(H_T)$ and that $\alpha T \ll 1$
\textcolor{green}{what does this mean and how does this relate to Prandtl number? 
and if $\alpha \ll 1$, how come $D$ is not $\ll 1$?} 
\subsection{Reference state and equation of state}
A reference state with hydrostatic pressure $P_H$ and adiabatic temperature distribution (gradient?) $T_s$ is considered. 
The equation of state is written as a linear Taylor expansion of $\rho$ about the reference state density distribution $\rho_r$: 
\begin{equation} \label{eq:11}
    \rho(T,P) = \rho_r[1-\alpha(T-T_s)+K_T^{-1}(P-P_H)]
\end{equation}
Where 
\begin{align}
    \alpha &= \frac{1}{\rho_r} \left(\pdv{\rho_R}{T}\right)_P \\
    K_T &= \rho_r \left(\pdv{P}{\rho_r}\right)_T 
\end{align}

Determining the reference density distribution $\rho_r$ is done by considering the gradient of $\rho_r(S, P_H)$ with $S$ 
the entropy per unit mass: 
\begin{equation}
    \nabla\rho_R(S,P_H) = \left(\pdv{\rho_r}{S}\right)_{P_H} \nabla S + \left(\pdv{\rho_r}{P_H}\right)_S\nabla P_H
\end{equation}
With $\nabla S = 0$ since the reference state is adiabatic, and with $\nabla P_H = -\rho_rg \bm{\hat{z}}$:
\begin{equation}
    \frac{1}{\rho_r} \pdv{\rho_r}{z} = -\frac{\rho_r g}{K_s}
\end{equation} 
with $K_s = \rho_r\left(\pdv{P_H}{\rho_r}\right)_s$. Using the definition of the Grüneisen parameter $\Gamma$ and the temperature scale height of the liquid
$H_T$ gives the Adams-Williamson relation \textcolor{green}{add images!}
\begin{equation} \label{eq:adams-williamson}
    \frac{1}{\rho_r} \pdv{\rho_r}{z} = -\frac{1}{H_T \Gamma}
\end{equation}
Which is valid for a \textbf{chemically homogeneous adiabatic density distribution under hydrostatic pressure}. 
Integration of equation \ref{eq:adams-williamson} over $z$ gives the reference densitiy distribution as
\begin{equation}
    \rho_r(z)= \rho_0 \exp[(d-z)/H_T \Gamma]
\end{equation}
where $\rho_0 = \rho_r(z=d)$ is the reference density value at the surface \textcolor{green}{orientation of coordinate system??}
Plugging this expression of the reference density distribution into \ref{eq:11} gives the (linearized equation) of state
\begin{equation} \label{eq:16}
    \rho(T,P) = \rho_0 \exp[(d-z)/H_T \Gamma][1-\alpha (T-T_s)+K_T^{-1}(P-P_H)]
\end{equation}

\subsection{Nondimensionalization}
\ref{table:characteristic_values}, \ref{table:characteristic_values}. \textcolor{green}{move tables here?}
In this section, \cite{jamc80} introduces various characteristic values of variables in order to nondimensionalize the preceding equations. These characteristic values and the resulting dimensionless 
variables are presented in the tables below.The characteristic value for \textbf{velocity} is taken directly from \cite{mcrw74}, and their derivation is expanded below: 

\begin{mdframed} 
    Our starting point is the momentum equation (in the notation of \cite{mcrw74}), with the left-hand side set to 0 (on the basis of $\mathrm{Re} \ll 1$ and $\frac{1}{\mathrm{Pr}}\to 0$):
    \begin{align} \label{eq:mcrw:1}
        0 &= -\frac{1}{\rho} \pdv{P}{x_i} + \pdv{\phi}{x_i} + \frac{1}{\rho}\pdv{\tau_{ij}}{x_j}
    \end{align}
    with 
    \begin{align*}
        \tau_{ij} &= \rho \nu\left(\pdv{u_i}{x_j} + \pdv{u_j}{x_i}\right) \\
        \nabla\phi &= \bm{g}
    \end{align*}
    last term in Eq.~\eqref{eq:mcrw:1} becomes \textcolor{green}{help index confusion}
    \begin{align*}
        \frac{1}{\rho}\pdv{\tau_{ij}}{x_j} &= \rho \nu \left(\pdv{x_j}\pdv{u_i}{x_j}+\pdv{x_j}\pdv{u_j}{x_i}\right) \\
        &= \rho \nu (\nabla^2u_i+ \pdv{x_i}\nabla \cdot \bm{u}) \\
    \end{align*}
    \textbf{assuming incompressibility,} $\nabla \cdot \bm{u} = 0$, and \textcolor{green}{or should this be $\bm{u}$ (vector)?}
    \begin{align} 
        \frac{1}{\rho}\pdv{\tau_{ij}}{x_j} = \nu \nabla^2 u_i
    \end{align}

    Now we formulate an equation of state under the Boussinesq approximation (Taylor expansion of density $\rho(T)$ around a reference state temperature $T_0$): 
    \begin{align*}
        \rho &= \rho_0[1-\alpha (T-T_0)] \\
        \frac{1}{\rho} &= \rho_0^{-1}[1-\alpha (T-T_0)]^{-1} \\
        &\approx \rho_0^{-1}[1+\alpha (T-T_0)]
    \end{align*}
    with $\alpha = -\frac{1}{\rho}\left(\pdv{\rho}{T}\right)_P$ the thermal expansivity coefficient. 

    We can split the pressure gradient $\nabla P$ into a hydrostatic part $(\nabla P_s)$ and a non-hydrostatic part $\nabla p$: 
    \begin{align*}
        \nabla P &= \nabla P_s + \nabla p  \\
        &= \rho_0 \bm{g} + \nabla p
    \end{align*}
    Together with the equation of state, the pressure term + gravity in Eq.~\eqref{eq:mcrw:1} then becomes 
    \begin{align*}
        -\frac{1}{\rho} \nabla P + \bm{g}&= -\rho_0^{-1}(1+\alpha(T-T_0))(\nabla P_s + \nabla p) + \bm{g}\\
        & = -\rho_0^{-1} \nabla P_s  - \rho_0^{-1}\nabla P_s \alpha(T-T_0) - \rho_0^{-1} \nabla p -  \rho_0^{-1} \nabla p \alpha(T-T_0) + \bm{g} \\
        & \approx -\bm{g}\alpha(T-T_0) - \rho_0^{-1}\nabla p
    \end{align*}
    Where we neglected the second-order terms ($\sim \nabla p T$) and used that $\nabla P_s = \rho_0\bm{g}$, cancelling out the gravity terms.
    With this, Eq.~\eqref{eq:mcrw:1} becomes 
    \begin{align} \label{eq:mcrw:10}
        \nu \nabla^2 \bm{u} = \bm{g}\alpha(T-T_0)+\rho_0^{-1}\nabla p
    \end{align} 
    \textbf{note that} $\nabla p$ \textbf{is only the non-hydrostatic pressure gradient}. 

    We can take the curl of Eq.~\eqref{eq:mcrw:10}, using that $\nabla \times \bm{u} = \bm{\omega}$, that the curl of a gradient is always zero ($\nabla \times (\nabla f) =0$), and that $\bm{g} = \nabla \phi$: 
    \begin{align} \label{eq:mcrw:13}
        \begin{split}
            \nu \nabla^2\bm{\omega} &= \nabla \times (\bm{g}\alpha(T-T_0))+\nabla \times (\rho_0^{-1}\nabla p) \\
        &= \nabla \times [\bm{g}\alpha(T-T_0)] \\
        &= \alpha(T-T_0)(\nabla \times \bm{g})+ \alpha \nabla T \times \bm{g} \\
        & = \alpha \nabla T \times \bm{g}
        \end{split}
    \end{align}
Now we assume that the flow in restricted to the $x,z$ plane, meaning 
$\bm{u} = (u,0,w) = \left(-\pdv{\psi}{z}, 0, \pdv{\psi}{x}\right)$ (where $\psi$ is the stream function) 
and thus $\bm{\omega} = (0,\omega, 0) = (0, -\nabla^2\psi, 0)$. With $\bm{g} = (0,0,g)$, Eq.~\ref{eq:mcrw:13} becomes (computing the cross product)
\begin{align} \label{eq:mcrw:16}
    \nabla^2\omega = \frac{\alpha g}{\nu}\pdv{T}{x}
\end{align}
From this equation, we can get an characteristic velocity for the convecting system. 
\begin{align*}
    \mathcal{O}(\nabla^2\omega) &=\mathcal{O}(\frac{\alpha g}{\nu} \pdv{T}{x}) \\
    \frac{U}{d^3} &= \frac{\alpha g}{\nu_0} \frac{T_1}{d} \\
    U &= \frac{\alpha g T_1 d^2}{\nu_0}
\end{align*}
Where $U$, $\nu_0$, and $T_1$ are characteristic values for velocity, kinematic viscosity, and temperature deviations from the adiabatic reference profile $T_0$, respectively. 
\end{mdframed}

For the scaling of pressure variations $P_1$ \textcolor{green}{TO DO}

\begin{table}[h!] 
    \renewcommand{\arraystretch}{1.8}   % taller rows
    \setlength{\tabcolsep}{12pt}        % wider column gaps
    \centering
    \caption{characteristic values of \cite{jamc80}}
    \label{table:nondimensionalisation}
    \begin{tabular}{l  >{$}l<{$} >{$}l<{$}}
        \toprule
        Variable & \text{Characteristic value} & \text{Dimensionless variable} \\
        \midrule
        depth & d & (x', z') = \dfrac{(x,z)}{d}\\
        temperature & T_1 (=1 \si{\kelvin}) & T' = \dfrac{(T-T_0)}{T_1}\\
        density & \rho_0 & \rho_r' = \dfrac{\rho_r}{\rho_0}, \quad \rho = \dfrac{\rho}{\rho_0} \\
        velocity & U = \dfrac{g\alpha T_1 d^2}{\nu _0} & \\
        kinematic viscosity & \nu_0 = \eta/ \rho_0 &  \\
        time & \tau = d/U & \\
        \bottomrule
    \end{tabular}
\end{table}




\subsection{The anelastic-liquid approximation}
From now on, only the dimensionless variables will be used but with primes surpressed (i.e. all variables are assumed dimensionless). 
The dimensionless equation of state (from eq. \ref{eq:16}) becomes 
\begin{equation}
    \rho = \rho_r[1-\mu(T-T_s)+\mu(\rho_0gd/K_T)P_1]
\end{equation}
with $\mu = \alpha T_1$. If $\mu \to 0$, $\rho$ reduces to its reference profile $\rho_r$ (as is the case in the Boussinesq approximation). 


\section{Appendix 1: tables of symbols}
\begin{table}[ht] 
    \renewcommand{\arraystretch}{1.8}   % taller rows
    \setlength{\tabcolsep}{12pt}        % wider column gaps
    \centering
    \caption{Symbols used by \cite{jamc80}}
    \label{table:symbols}
    \begin{tabular}{>{$}l<{$} >{$}l<{$} l}
        \toprule
        d & & depth of the convecting layer \\
        H_T & \dfrac{C_p}{g\alpha} & temperature scale height of the liquid \\
        C_p & & specific heat at constant pressure \\
        C_v && specific heat at constant volume \\
        \alpha & & volume coefficient of thermal expansion \\
        \tau_{ij} && deviatoric stress tensor \\
        (\nabla T)_s & \dfrac{\alpha T}{\rho C_p}\nabla P & adiabatic temperature gradient \\
        K && thermal conductivity \\
        H && rate of internal heat production \\
        \eta & \rho \nu & dynamic viscosity \\
        \nu && kinematic viscosity \\
        \eta_0 & \ll 1 & second or bulk viscosity \\
        \Gamma & \dfrac{\alpha K_s}{\rho C_p} = \dfrac{\alpha K_T}{\rho C_v} & Grüneisen's parameter, assumed constant and $O(1)$\\ 
        K_s & & isentropic modulus of bulk compressibility \\
        K_T && isothermal modulus of bulk compressibility \\ 
        \mu & \alpha T_1 & \\
        \alpha & 10^{-5} \quad ^\circ K^{-1} & volume coefficient of thermal expansion \\ 



        % T_0 & & reference temperature \\
        % \Delta T^* & & temperature increase relative to reference temperature \\
        \bottomrule
    \end{tabular}
\end{table}

\begin{table}[ht] 
    \renewcommand{\arraystretch}{1.8}   % taller rows
    \setlength{\tabcolsep}{12pt}        % wider column gaps
    \centering
    \caption{characteristic values of \cite{jamc80}}
    \label{table:characteristic_values}
    \begin{tabular}{l  >{$}l<{$}  }
        \toprule
        depth & d\\
        temperature & T_1 = \dfrac{1}{\text{K}}\textcolor{green}{WHAT??} \to 1 ^{\circ} \text{K}\\
        density & \rho_0 \\
        velocity & U = \dfrac{g\alpha T_1 d^2}{\nu _0}, \quad \nu_0 = \eta/ \rho_0 \text{\textcolor{green}{help this is different from schubert}} \\
        time & \tau = d/U \\
        \bottomrule
    \end{tabular}
\end{table}

\begin{table}[ht] 
    \renewcommand{\arraystretch}{1.8}   % taller rows
    \setlength{\tabcolsep}{12pt}        % wider column gaps
    \centering
    \caption{Dimensionless variables in \cite{jamc80}}
    \label{table:dimensionless_variables}
    \begin{tabular}{>{$}l<{$} >{$}l<{$} }
        \toprule
        (x', z') & (x,z)/d\\
        \rho_r'& \rho_r/\rho_0 = \exp[(1-z')D/\Gamma], \quad D= d/H_T\\
        \rho' & \rho/\rho_0 \\
        T'& (T-T_0)/T_1 \\ 
        t' & t/\tau \\
        P_1' & \dfrac{P_1 d}{U\eta} = P_1 \alpha T_1 \rho_0 gd, \quad P_1 = P-P_H \\
        \kappa'_0 & \kappa_0\nu_0/g\alpha T_1d^3, \quad \text{thermal diffusivity} \\
        \epsilon'_0 & \epsilon_0\nu_0/g\alpha T_1^2d, \quad \text{energy generation rate,}\quad \epsilon_0 = H/\rho_0C_p\\
        \lambda & L/d, \quad \text{box width, } \quad L \quad \text{is the dimensional box width}\\
        \bottomrule
    \end{tabular}
\end{table}
\end{document}